\chapter{Доказательства теоретических результатов}\label{app:transformer-details}

\section{Гаусс-ньютоновская компонента матрицы Гессе для трансформерной архитектуры нейронной сети}\label{app:transformer-details}

В данном приложении приводятся детальные выкладки, на которые опирается основной текст раздела о трансформерных архитектурах. Мы сохраняем обозначения основного текста: рассматривается одно-головый трансформерный блок архитектуры Post-LN
\begin{equation}
    \mathbf{U} = \mathbf{X} + \mathrm{SA}(\mathbf{X}),
    \qquad
    \mathbf{Y} = \mathrm{LN}(\mathbf{U}),
\end{equation}
\begin{equation}
    \mathrm{FFN}(\mathbf{Y}) = \sigma(\mathbf{Y}\mathbf{W}_1)\mathbf{W}_2,
    \qquad
    \mathbf{S} = \mathbf{Y} + \mathrm{FFN}(\mathbf{Y}),
    \qquad
    \mathbf{Z} = \mathrm{LN}(\mathbf{S}),
\end{equation}
где
\begin{equation}
    \mathrm{SA}(\mathbf{X})
    =
    \mathbf{\Pi}(\mathbf{X})\mathbf{X}\mathbf{W}_V,
    \qquad
    \mathbf{\Pi}(\mathbf{X})
    =
    \softmax\left(
        \frac{\mathbf{X}\mathbf{W}_Q\mathbf{W}_K^\top \mathbf{X}^\top}{\sqrt{d_k}}
    \right).
\end{equation}
Выход блока обозначается через
\begin{equation}
    \mathbf{z} = \vecrop(\mathbf{Z}) \in \mathbb{R}^{Td}.
\end{equation}

\subsection{Якобиан и оценка нормы оператора LayerNorm}

Рассмотрим произвольную матрицу
\begin{equation}
    \mathbf{U} \in \mathbb{R}^{T\times d}.
\end{equation}
Введем матрицу центрирования
\begin{equation}
    \mathbf{C}_d
    =
    \mathbf{I}_d - \frac{1}{d}\mathbf{1}_d\mathbf{1}_d^\top
    \in \mathbb{R}^{d\times d}.
\end{equation}
Тогда центрированное представление строки задается оператором
\begin{equation}
    \mathbf{M}(\mathbf{U})
    =
    \mathbf{U}\mathbf{C}_d
    =
    \mathbf{U} - \frac{1}{d}\mathbf{U}\mathbf{1}_d\mathbf{1}_d^\top.
\end{equation}
Вектор построчных стандартных отклонений определим формулой
\begin{equation}
    \boldsymbol{\sigma}(\mathbf{U})
    =
    \frac{1}{\sqrt d}
    \left(
        \mathbf{M}(\mathbf{U})^{\circ 2}\mathbf{1}_d
    \right)^{\circ \frac12}
    \in \mathbb{R}^{T},
\end{equation}
а диагональную матрицу обратных стандартных отклонений --- как
\begin{equation}
    \mathbf{P}(\mathbf{U})
    =
    \diag^{-1}\left(\boldsymbol{\sigma}(\mathbf{U})\right)
    \in \mathbb{R}^{T\times T}.
\end{equation}
Тогда LayerNorm без аффинных параметров $\gamma,\beta$ записывается в матричном виде:
\begin{equation}
    \label{eq:app-tr-LN-form}
    \mathrm{LN}(\mathbf{U})
    =
    \mathbf{P}(\mathbf{U})\mathbf{M}(\mathbf{U}).
\end{equation}

Для дальнейшего дифференцирования будем предполагать, что все построчные дисперсии положительны:
\begin{equation}
    \label{eq:app-tr-ln-positive-var}
    \sigma_{\min}(\mathbf{U})
    :=
    \min_{t=1,\ldots,T}\sigma_t(\mathbf{U}) > 0.
\end{equation}
Это стандартное техническое предположение, используемое и в статье-источнике при выводе якобиана LayerNorm. 

\begin{lemma}
    \label{lemma:app-tr-LN-jacobian}
    Якобиан LayerNorm по входу имеет вид
    \begin{align}
        \label{eq:app-tr-LN-jacobian}
        \mathbf{J}_{\mathrm{LN}}(\mathbf{U})
        &=
        \frac{\partial \vecrop(\mathrm{LN}(\mathbf{U}))}{\partial \vecrop(\mathbf{U})^\top}
        \notag\\
        &=
        \left(
            \mathbf{P}(\mathbf{U}) \otimes \mathbf{I}_d
        \right)\mathbf{G}_d
        +
        \left(
            \mathbf{I}_T \otimes \mathbf{M}(\mathbf{U})^\top
        \right)
        \frac{\partial \vecrop(\mathbf{P}(\mathbf{U}))}{\partial \vecrop(\mathbf{U})^\top},
    \end{align}
    где
    \begin{equation}
        \mathbf{G}_d
        =
        \frac{\partial \vecrop(\mathbf{M}(\mathbf{U}))}{\partial \vecrop(\mathbf{U})^\top}
        =
        \mathbf{I}_{Td}
        -
        \frac{1}{d}
        \left(
            \mathbf{I}_T \otimes \mathbf{1}_{d\times d}
        \right).
    \end{equation}
\end{lemma}

\begin{proof}
Из \eqref{eq:app-tr-LN-form} имеем
\begin{equation}
    \mathrm{LN}(\mathbf{U})
    =
    \mathbf{P}(\mathbf{U})\mathbf{M}(\mathbf{U}).
\end{equation}
Применим правило дифференцирования матричного произведения к построчной векторизации:
\begin{equation}
    \frac{\partial \vecrop(\mathbf{P}(\mathbf{U})\mathbf{M}(\mathbf{U}))}{\partial \vecrop(\mathbf{U})^\top}
    =
    \left(
        \mathbf{P}(\mathbf{U}) \otimes \mathbf{I}_d
    \right)
    \frac{\partial \vecrop(\mathbf{M}(\mathbf{U}))}{\partial \vecrop(\mathbf{U})^\top}
    +
    \left(
        \mathbf{I}_T \otimes \mathbf{M}(\mathbf{U})^\top
    \right)
    \frac{\partial \vecrop(\mathbf{P}(\mathbf{U}))}{\partial \vecrop(\mathbf{U})^\top}.
\end{equation}
Остается вычислить производную центрирующего оператора. Так как
\begin{equation}
    \mathbf{M}(\mathbf{U}) = \mathbf{U}\mathbf{C}_d,
\end{equation}
то для построчной векторизации выполняется
\begin{equation}
    \vecrop(\mathbf{M}(\mathbf{U}))
    =
    (\mathbf{I}_T \otimes \mathbf{C}_d)\vecrop(\mathbf{U}).
\end{equation}
Следовательно,
\begin{equation}
    \frac{\partial \vecrop(\mathbf{M}(\mathbf{U}))}{\partial \vecrop(\mathbf{U})^\top}
    =
    \mathbf{I}_T \otimes \mathbf{C}_d
    =
    \mathbf{I}_{Td}
    -
    \frac{1}{d}
    \left(
        \mathbf{I}_T \otimes \mathbf{1}_{d\times d}
    \right)
    =
    \mathbf{G}_d.
\end{equation}
Подставляя это равенство в предыдущую формулу, получаем \eqref{eq:app-tr-LN-jacobian}. Данная форма совпадает с матричным представлением якобиана LayerNorm в статье-источнике. 
\end{proof}

\begin{lemma}
    \label{lemma:app-tr-LN-jacobian-bound}
    Пусть выполнено условие \eqref{eq:app-tr-ln-positive-var}. Тогда
    \begin{equation}
        \label{eq:app-tr-LN-bound}
        \left\|
            \mathbf{J}_{\mathrm{LN}}(\mathbf{U})
        \right\|_2
        \le
        \frac{1}{\sigma_{\min}(\mathbf{U})}
        +
        \frac{\|\mathbf{U}\|_2^2}{\sqrt{d}\,\sigma_{\min}(\mathbf{U})^3}.
    \end{equation}
\end{lemma}

\begin{proof}
По лемме \ref{lemma:app-tr-LN-jacobian}
\begin{equation}
    \mathbf{J}_{\mathrm{LN}}(\mathbf{U})
    =
    \left(
        \mathbf{P}(\mathbf{U}) \otimes \mathbf{I}_d
    \right)\mathbf{G}_d
    +
    \left(
        \mathbf{I}_T \otimes \mathbf{M}(\mathbf{U})^\top
    \right)
    \frac{\partial \vecrop(\mathbf{P}(\mathbf{U}))}{\partial \vecrop(\mathbf{U})^\top}.
\end{equation}
Следовательно, по неравенству треугольника и субмультипликативности спектральной нормы
\begin{align}
    \left\|
        \mathbf{J}_{\mathrm{LN}}(\mathbf{U})
    \right\|_2
    &\le
    \left\|
        \mathbf{P}(\mathbf{U}) \otimes \mathbf{I}_d
    \right\|_2
    \|\mathbf{G}_d\|_2
    +
    \left\|
        \mathbf{I}_T \otimes \mathbf{M}(\mathbf{U})^\top
    \right\|_2
    \left\|
        \frac{\partial \vecrop(\mathbf{P}(\mathbf{U}))}{\partial \vecrop(\mathbf{U})^\top}
    \right\|_2.
\end{align}
Используя свойство нормы произведения Кронекера, получаем
\begin{equation}
    \left\|
        \mathbf{P}(\mathbf{U}) \otimes \mathbf{I}_d
    \right\|_2
    =
    \|\mathbf{P}(\mathbf{U})\|_2,
    \qquad
    \left\|
        \mathbf{I}_T \otimes \mathbf{M}(\mathbf{U})^\top
    \right\|_2
    =
    \|\mathbf{M}(\mathbf{U})\|_2.
\end{equation}
Поэтому
\begin{equation}
    \label{eq:app-tr-ln-bound-step-1}
    \left\|
        \mathbf{J}_{\mathrm{LN}}(\mathbf{U})
    \right\|_2
    \le
    \|\mathbf{P}(\mathbf{U})\|_2 \|\mathbf{G}_d\|_2
    +
    \|\mathbf{M}(\mathbf{U})\|_2
    \left\|
        \frac{\partial \vecrop(\mathbf{P}(\mathbf{U}))}{\partial \vecrop(\mathbf{U})^\top}
    \right\|_2.
\end{equation}




\textbf{Шаг 1. Оценка нормы $\mathbf{G}_d$.}
Так как
\begin{equation}
    \mathbf{G}_d = \mathbf{I}_T \otimes \mathbf{C}_d,
\end{equation}
а матрица $\mathbf{C}_d$ есть ортогональный проектор на подпространство векторов с нулевой суммой координат, то
\begin{equation}
    \|\mathbf{C}_d\|_2 = 1,
    \qquad
    \|\mathbf{G}_d\|_2
    =
    \|\mathbf{I}_T \otimes \mathbf{C}_d\|_2
    =
    1.
\end{equation}




\textbf{Шаг 2. Оценка нормы $\mathbf{P}(\mathbf{U})$.}
Матрица $\mathbf{P}(\mathbf{U})$ диагональна, и на ее диагонали стоят числа
\begin{equation}
    \frac{1}{\sigma_1(\mathbf{U})}, \ldots, \frac{1}{\sigma_T(\mathbf{U})}.
\end{equation}
Следовательно,
\begin{equation}
    \|\mathbf{P}(\mathbf{U})\|_2
    =
    \max_{t=1,\ldots,T}\frac{1}{\sigma_t(\mathbf{U})}
    =
    \frac{1}{\sigma_{\min}(\mathbf{U})}.
\end{equation}




\textbf{Шаг 3. Оценка нормы $\mathbf{M}(\mathbf{U})$.}
Поскольку
\begin{equation}
    \mathbf{M}(\mathbf{U}) = \mathbf{U}\mathbf{C}_d,
\end{equation}
имеем
\begin{equation}
    \|\mathbf{M}(\mathbf{U})\|_2
    \le
    \|\mathbf{U}\|_2 \|\mathbf{C}_d\|_2
    \le
    \|\mathbf{U}\|_2.
\end{equation}




\textbf{Шаг 4. Оценка нормы производной $\partial \vecrop(\mathbf{P}(\mathbf{U}))/\partial \vecrop(\mathbf{U})^\top$.}
Рассмотрим фиксированную строку $\mathbf{m}_t^\top \in \mathbb{R}^{1\times d}$ матрицы $\mathbf{M}(\mathbf{U})$. Тогда
\begin{equation}
    \sigma_t(\mathbf{U})
    =
    \frac{1}{\sqrt d}\|\mathbf{m}_t\|_2.
\end{equation}
Следовательно,
\begin{equation}
    \frac{1}{\sigma_t(\mathbf{U})}
    =
    \frac{\sqrt d}{\|\mathbf{m}_t\|_2},
\end{equation}
и потому
\begin{equation}
    \frac{\partial}{\partial \mathbf{m}_t}
    \left(
        \frac{1}{\sigma_t(\mathbf{U})}
    \right)
    =
    -\sqrt d\,
    \frac{\mathbf{m}_t^\top}{\|\mathbf{m}_t\|_2^3}.
\end{equation}
Отсюда
\begin{equation}
    \left\|
        \frac{\partial}{\partial \mathbf{m}_t}
        \left(
            \frac{1}{\sigma_t(\mathbf{U})}
        \right)
    \right\|_2
    =
    \sqrt d \frac{1}{\|\mathbf{m}_t\|_2^2}
    =
    \frac{1}{\sqrt d\,\sigma_t(\mathbf{U})^2}.
\end{equation}
Так как отображение $\mathbf{U} \mapsto \mathbf{M}(\mathbf{U})=\mathbf{U}\mathbf{C}_d$ линейно и имеет операторную норму не больше единицы, то, переходя от производной по $\mathbf{m}_t$ к производной по $\mathbf{U}$, получаем
\begin{equation}
    \left\|
        \frac{\partial}{\partial \vecrop(\mathbf{U})}
        \left(
            \frac{1}{\sigma_t(\mathbf{U})}
        \right)
    \right\|_2
    \le
    \frac{1}{\sqrt d\,\sigma_t(\mathbf{U})^2}.
\end{equation}
Поскольку $\mathbf{P}(\mathbf{U})$ есть диагональная матрица, составленная из величин $1/\sigma_t(\mathbf{U})$, то оператор дифференцирования диагонального отображения не увеличивает порядок оценки, и потому
\begin{equation}
    \left\|
        \frac{\partial \vecrop(\mathbf{P}(\mathbf{U}))}{\partial \vecrop(\mathbf{U})^\top}
    \right\|_2
    \le
    \frac{1}{\sqrt d\,\sigma_{\min}(\mathbf{U})^2}.
\end{equation}
Для согласования с грубой глобальной оценкой, используемой в основном тексте, умножим правую часть на $\|\mathbf{U}\|_2/\sigma_{\min}(\mathbf{U})$, что дает более слабую, но удобную в дальнейшем форму:
\begin{equation}
    \left\|
        \frac{\partial \vecrop(\mathbf{P}(\mathbf{U}))}{\partial \vecrop(\mathbf{U})^\top}
    \right\|_2
    \le
    \frac{\|\mathbf{U}\|_2}{\sqrt d\,\sigma_{\min}(\mathbf{U})^3}.
\end{equation}



Подставляя все полученные оценки в \eqref{eq:app-tr-ln-bound-step-1}, получаем
\begin{equation}
    \left\|
        \mathbf{J}_{\mathrm{LN}}(\mathbf{U})
    \right\|_2
    \le
    \frac{1}{\sigma_{\min}(\mathbf{U})}
    +
    \|\mathbf{U}\|_2
    \cdot
    \frac{\|\mathbf{U}\|_2}{\sqrt d\,\sigma_{\min}(\mathbf{U})^3},
\end{equation}
то есть
\begin{equation}
    \left\|
        \mathbf{J}_{\mathrm{LN}}(\mathbf{U})
    \right\|_2
    \le
    \frac{1}{\sigma_{\min}(\mathbf{U})}
    +
    \frac{\|\mathbf{U}\|_2^2}{\sqrt d\,\sigma_{\min}(\mathbf{U})^3}.
\end{equation}
Что и требовалось доказать. Такая же структура оценки используется в статье для контроля нормы якобиана LayerNorm. 
\end{proof}

\subsection{Выкладки для feed-forward части}

Далее полагаем $\sigma = \mathrm{ReLU}$, действующую покомпонентно. Введем обозначение
\begin{equation}
    \mathbf{D}_{\sigma}(\mathbf{Y}\mathbf{W}_1)
    =
    \diag\left(
        \vecrop\left(
            \mathbf{1}_{\{\mathbf{Y}\mathbf{W}_1 > 0\}}
        \right)
    \right)
    \in \mathbb{R}^{Td_{\mathrm{ff}}\times Td_{\mathrm{ff}}}.
\end{equation}

\begin{lemma}
    \label{lemma:app-tr-ReLU}
    Почти всюду по $\mathbf{Y}\mathbf{W}_1$ справедливо
    \begin{equation}
        \frac{\partial \vecrop(\sigma(\mathbf{Y}\mathbf{W}_1))}{\partial \vecrop(\mathbf{Y}\mathbf{W}_1)^\top}
        =
        \mathbf{D}_{\sigma}(\mathbf{Y}\mathbf{W}_1),
    \end{equation}
    причем вторая производная ReLU почти всюду равна нулю.
\end{lemma}

\begin{proof}
ReLU действует покомпонентно, поэтому по каждой координате ее производная равна $1$ на положительной полуоси и $0$ на отрицательной. Отсюда якобиан покомпонентного отображения есть диагональная матрица, на диагонали которой стоят соответствующие индикаторы положительности. В точках, где аргумент равен нулю, производная не определена, но это множество имеет меру нуль. Вне этого множества якобиан кусочно постоянен, следовательно его производная почти всюду равна нулю. Тот же результат используется в статье как Lemma 4. 
\end{proof}

\begin{proposition}
    \label{prop:app-tr-FFN-jacobians}
    Для отображения
    \begin{equation}
        \mathbf{F}(\mathbf{Y})
        =
        \sigma(\mathbf{Y}\mathbf{W}_1)\mathbf{W}_2
        \in \mathbb{R}^{T\times d}
    \end{equation}
    параметрические якобианы по матрицам $\mathbf{W}_1$ и $\mathbf{W}_2$ задаются формулами
    \begin{equation}
        \label{eq:app-tr-FFN-jac-W2}
        \frac{\partial \vecrop(\mathbf{F}(\mathbf{Y}))}{\partial \mathbf{w}_2^\top}
        =
        \sigma(\mathbf{Y}\mathbf{W}_1)\otimes \mathbf{I}_d,
    \end{equation}
    \begin{equation}
        \label{eq:app-tr-FFN-jac-W1}
        \frac{\partial \vecrop(\mathbf{F}(\mathbf{Y}))}{\partial \mathbf{w}_1^\top}
        =
        \left(
            \mathbf{I}_T \otimes \mathbf{W}_2^\top
        \right)
        \mathbf{D}_{\sigma}(\mathbf{Y}\mathbf{W}_1)
        \left(
            \mathbf{Y} \otimes \mathbf{I}_{d_{\mathrm{ff}}}
        \right).
    \end{equation}
\end{proposition}

\begin{proof}
Положим
\begin{equation}
    \mathbf{H} = \sigma(\mathbf{Y}\mathbf{W}_1)\in\mathbb{R}^{T\times d_{\mathrm{ff}}}.
\end{equation}
Тогда
\begin{equation}
    \mathbf{F}(\mathbf{Y}) = \mathbf{H}\mathbf{W}_2.
\end{equation}




\textbf{Производная по $\mathbf{W}_2$.}
Для построчной векторизации выполняется тождество
\begin{equation}
    \vecrop(\mathbf{H}\mathbf{W}_2)
    =
    (\mathbf{H}\otimes \mathbf{I}_d)\vecrop(\mathbf{W}_2).
\end{equation}
Следовательно,
\begin{equation}
    \frac{\partial \vecrop(\mathbf{F}(\mathbf{Y}))}{\partial \mathbf{w}_2^\top}
    =
    \mathbf{H}\otimes \mathbf{I}_d
    =
    \sigma(\mathbf{Y}\mathbf{W}_1)\otimes \mathbf{I}_d.
\end{equation}
Это и есть формула \eqref{eq:app-tr-FFN-jac-W2}.




\textbf{Производная по $\mathbf{W}_1$.}
Рассмотрим композицию
\begin{equation}
    \mathbf{W}_1
    \longmapsto
    \mathbf{Y}\mathbf{W}_1
    \longmapsto
    \sigma(\mathbf{Y}\mathbf{W}_1)
    \longmapsto
    \sigma(\mathbf{Y}\mathbf{W}_1)\mathbf{W}_2.
\end{equation}
По правилу дифференцирования сложной функции
\begin{equation}
    \frac{\partial \vecrop(\mathbf{F}(\mathbf{Y}))}{\partial \mathbf{w}_1^\top}
    =
    \frac{\partial \vecrop(\sigma(\mathbf{Y}\mathbf{W}_1)\mathbf{W}_2)}{\partial \vecrop(\sigma(\mathbf{Y}\mathbf{W}_1))^\top}
    \cdot
    \frac{\partial \vecrop(\sigma(\mathbf{Y}\mathbf{W}_1))}{\partial \vecrop(\mathbf{Y}\mathbf{W}_1)^\top}
    \cdot
    \frac{\partial \vecrop(\mathbf{Y}\mathbf{W}_1)}{\partial \mathbf{w}_1^\top}.
\end{equation}
Первый множитель равен
\begin{equation}
    \mathbf{I}_T \otimes \mathbf{W}_2^\top.
\end{equation}
Второй по лемме \ref{lemma:app-tr-ReLU} имеет вид
\begin{equation}
    \mathbf{D}_{\sigma}(\mathbf{Y}\mathbf{W}_1).
\end{equation}
Для третьего множителя используем тождество
\begin{equation}
    \vecrop(\mathbf{Y}\mathbf{W}_1)
    =
    (\mathbf{Y}\otimes \mathbf{I}_{d_{\mathrm{ff}}})\vecrop(\mathbf{W}_1),
\end{equation}
откуда
\begin{equation}
    \frac{\partial \vecrop(\mathbf{Y}\mathbf{W}_1)}{\partial \mathbf{w}_1^\top}
    =
    \mathbf{Y}\otimes \mathbf{I}_{d_{\mathrm{ff}}}.
\end{equation}
Перемножая эти три матрицы, получаем \eqref{eq:app-tr-FFN-jac-W1}. Именно эти формулы используются и в статье как FFN Jacobians. 
\end{proof}

\begin{lemma}
    \label{lemma:app-tr-FFN-input-jacobian}
    Якобиан $\vecrop(\mathrm{FFN}(\mathbf{Y}))$ по входу $\vecrop(\mathbf{Y})$ имеет вид
    \begin{equation}
        \label{eq:app-tr-FFN-input-jacobian}
        \mathbf{J}_{\mathrm{FFN},Y}
        =
        \frac{\partial \vecrop(\mathrm{FFN}(\mathbf{Y}))}{\partial \vecrop(\mathbf{Y})^\top}
        =
        \left(
            \mathbf{I}_T \otimes \mathbf{W}_2^\top
        \right)
        \mathbf{D}_{\sigma}(\mathbf{Y}\mathbf{W}_1)
        \left(
            \mathbf{I}_T \otimes \mathbf{W}_1^\top
        \right).
    \end{equation}
\end{lemma}

\begin{proof}
Композиция
\[
\mathbf{Y}
\longmapsto
\mathbf{Y}\mathbf{W}_1
\longmapsto
\sigma(\mathbf{Y}\mathbf{W}_1)
\longmapsto
\sigma(\mathbf{Y}\mathbf{W}_1)\mathbf{W}_2
\]
дифференцируется так же, как и в доказательстве предложения \ref{prop:app-tr-FFN-jacobians}. Отличие состоит лишь в последнем множителе:
\begin{equation}
    \frac{\partial \vecrop(\mathbf{Y}\mathbf{W}_1)}{\partial \vecrop(\mathbf{Y})^\top}
    =
    \mathbf{I}_T \otimes \mathbf{W}_1^\top.
\end{equation}
Подставляя этот множитель вместо $\mathbf{Y}\otimes \mathbf{I}_{d_{\mathrm{ff}}}$, получаем формулу \eqref{eq:app-tr-FFN-input-jacobian}.
\end{proof}

\subsection{Оценки для self-attention части}

Здесь мы используем обозначения
\begin{equation}
    \mathbf{G}_Q
    =
    \frac{\partial \vecrop(\mathrm{SA}(\mathbf{X}))}{\partial \mathbf{w}_Q^\top},
    \qquad
    \mathbf{G}_K
    =
    \frac{\partial \vecrop(\mathrm{SA}(\mathbf{X}))}{\partial \mathbf{w}_K^\top},
    \qquad
    \mathbf{G}_V
    =
    \frac{\partial \vecrop(\mathrm{SA}(\mathbf{X}))}{\partial \mathbf{w}_V^\top}.
\end{equation}

Для матрицы $\mathbf{G}_V$ вывод можно сделать непосредственно. Для матриц $\mathbf{G}_Q$ и $\mathbf{G}_K$ в данной главе мы опираемся на известный анализ изолированного single-head self-attention блока~\todo{(уточнить откуда, ссылка)}: именно эти оценки затем подставляются в якобиан полного трансформерного блока.

\begin{lemma}
    \label{lemma:app-tr-SA-jacobian-bounds}
    Для матриц $\mathbf{G}_Q,\mathbf{G}_K,\mathbf{G}_V$ справедливы оценки
    \begin{equation}
        \label{eq:app-tr-GV-bound}
        \|\mathbf{G}_V\|_2
        \le
        T\|\mathbf{X}\|_2,
    \end{equation}
    а также
    \begin{equation}
        \label{eq:app-tr-GQ-bound}
        \|\mathbf{G}_Q\|_2
        \le
        \frac{C_{\mathrm{sm}}}{\sqrt{d_k}}
        \|\mathbf{W}_V\|_2
        \|\mathbf{W}_K\|_2
        \|\mathbf{X}\|_2^3,
    \end{equation}
    \begin{equation}
        \label{eq:app-tr-GK-bound}
        \|\mathbf{G}_K\|_2
        \le
        \frac{C_{\mathrm{sm}}}{\sqrt{d_k}}
        \|\mathbf{W}_V\|_2
        \|\mathbf{W}_Q\|_2
        \|\mathbf{X}\|_2^3,
    \end{equation}
    где $C_{\mathrm{sm}}$ --- константа Липшица построчного softmax в спектральной норме.
\end{lemma}

\begin{proof}
\textbf{1. Оценка для $\mathbf{G}_V$.}
При фиксированных $\mathbf{X},\mathbf{W}_Q,\mathbf{W}_K$ оператор self-attention линеен по $\mathbf{W}_V$:
\begin{equation}
    \mathrm{SA}(\mathbf{X})
    =
    \mathbf{\Pi}(\mathbf{X})\mathbf{X}\mathbf{W}_V.
\end{equation}
Следовательно,
\begin{equation}
    \vecrop(\mathrm{SA}(\mathbf{X}))
    =
    \left(
        \mathbf{\Pi}(\mathbf{X})\mathbf{X}\otimes \mathbf{I}_d
    \right)\vecrop(\mathbf{W}_V),
\end{equation}
и потому
\begin{equation}
    \mathbf{G}_V
    =
    \mathbf{\Pi}(\mathbf{X})\mathbf{X}\otimes \mathbf{I}_d.
\end{equation}
По свойству спектральной нормы произведения Кронекера
\begin{equation}
    \|\mathbf{G}_V\|_2
    =
    \|\mathbf{\Pi}(\mathbf{X})\mathbf{X}\|_2
    \le
    \|\mathbf{\Pi}(\mathbf{X})\|_2 \|\mathbf{X}\|_2.
\end{equation}
Матрица $\mathbf{\Pi}(\mathbf{X})$ построчно стохастична, а потому грубо
\begin{equation}
    \|\mathbf{\Pi}(\mathbf{X})\|_2 \le T.
\end{equation}
Следовательно,
\begin{equation}
    \|\mathbf{G}_V\|_2 \le T\|\mathbf{X}\|_2.
\end{equation}



\textbf{2. Оценки для $\mathbf{G}_Q$ и $\mathbf{G}_K$.}
Рассмотрим матрицу логитов
\begin{equation}
    \mathbf{L}(\mathbf{X})
    =
    \frac{\mathbf{X}\mathbf{W}_Q\mathbf{W}_K^\top\mathbf{X}^\top}{\sqrt{d_k}}.
\end{equation}
При вариации $d\mathbf{W}_Q$ имеем
\begin{equation}
    d\mathbf{L}
    =
    \frac{\mathbf{X}\,d\mathbf{W}_Q\,\mathbf{W}_K^\top \mathbf{X}^\top}{\sqrt{d_k}},
\end{equation}
и потому
\begin{equation}
    \|d\mathbf{L}\|_2
    \le
    \frac{1}{\sqrt{d_k}}
    \|\mathbf{X}\|_2^2
    \|\mathbf{W}_K\|_2
    \|d\mathbf{W}_Q\|_2.
\end{equation}
Поскольку построчный softmax липшицев в спектральной норме с константой $C_{\mathrm{sm}}$, получаем
\begin{equation}
    \|d\mathbf{\Pi}\|_2
    \le
    C_{\mathrm{sm}}\|d\mathbf{L}\|_2
    \le
    \frac{C_{\mathrm{sm}}}{\sqrt{d_k}}
    \|\mathbf{X}\|_2^2
    \|\mathbf{W}_K\|_2
    \|d\mathbf{W}_Q\|_2.
\end{equation}
Далее
\begin{equation}
    d\,\mathrm{SA}(\mathbf{X})
    =
    (d\mathbf{\Pi})\mathbf{X}\mathbf{W}_V,
\end{equation}
поэтому
\begin{align}
    \|d\,\mathrm{SA}(\mathbf{X})\|_2
    &\le
    \|d\mathbf{\Pi}\|_2
    \|\mathbf{X}\mathbf{W}_V\|_2 \\
    &\le
    \frac{C_{\mathrm{sm}}}{\sqrt{d_k}}
    \|\mathbf{W}_K\|_2
    \|\mathbf{X}\|_2^2
    \|d\mathbf{W}_Q\|_2
    \|\mathbf{X}\|_2
    \|\mathbf{W}_V\|_2 \\
    &=
    \frac{C_{\mathrm{sm}}}{\sqrt{d_k}}
    \|\mathbf{W}_V\|_2
    \|\mathbf{W}_K\|_2
    \|\mathbf{X}\|_2^3
    \|d\mathbf{W}_Q\|_2.
\end{align}
Из определения операторной нормы якобиана следует \eqref{eq:app-tr-GQ-bound}. Формула \eqref{eq:app-tr-GK-bound} выводится симметрично при вариации по $\mathbf{W}_K$.
\end{proof}

\subsection{Сборка якобианов полного трансформерного блока}

Введем обозначения
\begin{equation}
    \mathbf{J}_Y
    =
    \frac{\partial \vecrop(\mathbf{Y})}{\partial \vecrop(\mathbf{U})^\top}
    =
    \mathbf{J}_{\mathrm{LN}}(\mathbf{U}),
    \qquad
    \mathbf{J}_Z
    =
    \frac{\partial \vecrop(\mathbf{Z})}{\partial \vecrop(\mathbf{S})^\top}
    =
    \mathbf{J}_{\mathrm{LN}}(\mathbf{S}).
\end{equation}

\begin{theorem}
    \label{thm:app-tr-block-jacobians}
    Якобианы выхода трансформерного блока по всем пяти группам параметров задаются формулами
    \begin{equation}
        \mathbf{J}^{(1)}
        =
        \mathbf{J}_Z \mathbf{B}_1,
        \qquad
        \mathbf{B}_1
        =
        \left(
            \mathbf{I}_T \otimes \mathbf{W}_2^\top
        \right)
        \mathbf{D}_{\sigma}(\mathbf{Y}\mathbf{W}_1)
        \left(
            \mathbf{Y} \otimes \mathbf{I}_{d_{\mathrm{ff}}}
        \right),
    \end{equation}
    \begin{equation}
        \mathbf{J}^{(2)}
        =
        \mathbf{J}_Z \mathbf{B}_2,
        \qquad
        \mathbf{B}_2
        =
        \sigma(\mathbf{Y}\mathbf{W}_1)\otimes \mathbf{I}_d,
    \end{equation}
    \begin{equation}
        \mathbf{J}^{(Q)}
        =
        \mathbf{J}_Z \mathbf{J}_{SY}\mathbf{J}_Y \mathbf{G}_Q,
        \qquad
        \mathbf{J}^{(K)}
        =
        \mathbf{J}_Z \mathbf{J}_{SY}\mathbf{J}_Y \mathbf{G}_K,
        \qquad
        \mathbf{J}^{(V)}
        =
        \mathbf{J}_Z \mathbf{J}_{SY}\mathbf{J}_Y \mathbf{G}_V,
    \end{equation}
    где
    \begin{equation}
        \mathbf{J}_{SY}
        =
        \frac{\partial \vecrop(\mathbf{S})}{\partial \vecrop(\mathbf{Y})^\top}
        =
        \mathbf{I}_{Td}
        +
        \left(
            \mathbf{I}_T \otimes \mathbf{W}_2^\top
        \right)
        \mathbf{D}_{\sigma}(\mathbf{Y}\mathbf{W}_1)
        \left(
            \mathbf{I}_T \otimes \mathbf{W}_1^\top
        \right).
    \end{equation}
\end{theorem}

\begin{proof}
\textbf{1. Якобианы по $\mathbf{W}_1$ и $\mathbf{W}_2$.}
Из определения блока имеем
\begin{equation}
    \mathbf{Z} = \mathrm{LN}(\mathbf{S}),
    \qquad
    \mathbf{S} = \mathbf{Y} + \mathrm{FFN}(\mathbf{Y}).
\end{equation}
Параметры $\mathbf{W}_1,\mathbf{W}_2$ входят только через FFN, а матрица $\mathbf{Y}$ от них не зависит. Поэтому
\begin{equation}
    \mathbf{J}^{(1)}
    =
    \frac{\partial \vecrop(\mathbf{Z})}{\partial \mathbf{w}_1^\top}
    =
    \frac{\partial \vecrop(\mathbf{Z})}{\partial \vecrop(\mathbf{S})^\top}
    \cdot
    \frac{\partial \vecrop(\mathbf{S})}{\partial \mathbf{w}_1^\top}
    =
    \mathbf{J}_Z
    \frac{\partial \vecrop(\mathrm{FFN}(\mathbf{Y}))}{\partial \mathbf{w}_1^\top}.
\end{equation}
По предложению \ref{prop:app-tr-FFN-jacobians}
\begin{equation}
    \frac{\partial \vecrop(\mathrm{FFN}(\mathbf{Y}))}{\partial \mathbf{w}_1^\top}
    =
    \left(
        \mathbf{I}_T \otimes \mathbf{W}_2^\top
    \right)
    \mathbf{D}_{\sigma}(\mathbf{Y}\mathbf{W}_1)
    \left(
        \mathbf{Y}\otimes \mathbf{I}_{d_{\mathrm{ff}}}
    \right).
\end{equation}
Тем самым получаем формулу для $\mathbf{J}^{(1)}$. Аналогично
\begin{equation}
    \mathbf{J}^{(2)}
    =
    \frac{\partial \vecrop(\mathbf{Z})}{\partial \mathbf{w}_2^\top}
    =
    \mathbf{J}_Z
    \frac{\partial \vecrop(\mathrm{FFN}(\mathbf{Y}))}{\partial \mathbf{w}_2^\top}
    =
    \mathbf{J}_Z
    \left(
        \sigma(\mathbf{Y}\mathbf{W}_1)\otimes \mathbf{I}_d
    \right).
\end{equation}



\textbf{2. Якобианы по attention-параметрам.}
Пусть $k\in\{Q,K,V\}$. Тогда влияние параметра $\mathbf{w}_k$ на выход проходит по цепочке
\begin{equation}
    \mathbf{w}_k
    \longmapsto
    \mathrm{SA}(\mathbf{X})
    \longmapsto
    \mathbf{U}
    \longmapsto
    \mathbf{Y}
    \longmapsto
    \mathbf{S}
    \longmapsto
    \mathbf{Z}.
\end{equation}
Следовательно, по правилу дифференцирования сложной функции
\begin{equation}
    \mathbf{J}^{(k)}
    =
    \frac{\partial \vecrop(\mathbf{Z})}{\partial \mathbf{w}_k^\top}
    =
    \frac{\partial \vecrop(\mathbf{Z})}{\partial \vecrop(\mathbf{S})^\top}
    \cdot
    \frac{\partial \vecrop(\mathbf{S})}{\partial \vecrop(\mathbf{Y})^\top}
    \cdot
    \frac{\partial \vecrop(\mathbf{Y})}{\partial \vecrop(\mathbf{U})^\top}
    \cdot
    \frac{\partial \vecrop(\mathrm{SA}(\mathbf{X}))}{\partial \mathbf{w}_k^\top}.
\end{equation}
Первый множитель равен $\mathbf{J}_Z$, третий --- $\mathbf{J}_Y$, четвертый --- $\mathbf{G}_k$. Остается вычислить
\begin{equation}
    \mathbf{J}_{SY}
    =
    \frac{\partial \vecrop(\mathbf{S})}{\partial \vecrop(\mathbf{Y})^\top}.
\end{equation}
Так как
\begin{equation}
    \mathbf{S} = \mathbf{Y} + \mathrm{FFN}(\mathbf{Y}),
\end{equation}
то
\begin{equation}
    \mathbf{J}_{SY}
    =
    \mathbf{I}_{Td}
    +
    \frac{\partial \vecrop(\mathrm{FFN}(\mathbf{Y}))}{\partial \vecrop(\mathbf{Y})^\top}.
\end{equation}
По лемме \ref{lemma:app-tr-FFN-input-jacobian}
\begin{equation}
    \frac{\partial \vecrop(\mathrm{FFN}(\mathbf{Y}))}{\partial \vecrop(\mathbf{Y})^\top}
    =
    \left(
        \mathbf{I}_T \otimes \mathbf{W}_2^\top
    \right)
    \mathbf{D}_{\sigma}(\mathbf{Y}\mathbf{W}_1)
    \left(
        \mathbf{I}_T \otimes \mathbf{W}_1^\top
    \right).
\end{equation}
Подстановка завершает доказательство.
\end{proof}

\subsection{Гаусс--ньютоновская компонента и ее спектральная оценка}

Пусть
\begin{equation}
    \mathbf{A}(\mathbf{w})
    =
    \nabla_{\mathbf{z}}^2 \ell(\mathbf{z},\mathbf{y})
    \in \mathbb{R}^{Td\times Td}
\end{equation}
--- матрица кривизны функции потерь по выходу блока, а
\begin{equation}
    \mathbf{J}(\mathbf{w})
    =
    \left[
        \mathbf{J}^{(Q)},
        \mathbf{J}^{(K)},
        \mathbf{J}^{(V)},
        \mathbf{J}^{(1)},
        \mathbf{J}^{(2)}
    \right]
\end{equation}
--- полный якобиан выхода блока по всем параметрам.

\begin{theorem}
    \label{thm:app-tr-GN-general-bound}
    Пусть
    \begin{equation}
        \|\mathbf{A}(\mathbf{w})\|_2 \le M_{\mathbf{A}}.
    \end{equation}
    Тогда
    \begin{equation}
        \label{eq:app-tr-GN-general-bound}
        \|\mathbf{G}(\mathbf{w})\|_2
        \le
        M_{\mathbf{A}}
        \left(
            \|\mathbf{J}^{(Q)}\|_2^2
            +
            \|\mathbf{J}^{(K)}\|_2^2
            +
            \|\mathbf{J}^{(V)}\|_2^2
            +
            \|\mathbf{J}^{(1)}\|_2^2
            +
            \|\mathbf{J}^{(2)}\|_2^2
        \right),
    \end{equation}
    где
    \begin{equation}
        \mathbf{G}(\mathbf{w})
        =
        \mathbf{J}(\mathbf{w})^\top
        \mathbf{A}(\mathbf{w})
        \mathbf{J}(\mathbf{w}).
    \end{equation}
    Кроме того,
    \begin{equation}
        \|\mathbf{J}^{(1)}\|_2
        \le
        \|\mathbf{J}_Z\|_2
        \|\mathbf{W}_2\|_2
        \|\mathbf{Y}\|_2,
    \end{equation}
    \begin{equation}
        \|\mathbf{J}^{(2)}\|_2
        \le
        \|\mathbf{J}_Z\|_2
        \|\mathbf{Y}\|_2
        \|\mathbf{W}_1\|_2,
    \end{equation}
    \begin{equation}
        \|\mathbf{J}^{(k)}\|_2
        \le
        \|\mathbf{J}_Z\|_2
        \left(
            \|\mathbf{W}_2\|_2\|\mathbf{W}_1\|_2 + 1
        \right)
        \|\mathbf{J}_Y\|_2
        \|\mathbf{G}_k\|_2,
        \qquad
        k\in\{Q,K,V\}.
    \end{equation}
\end{theorem}

\begin{proof}
По определению гаусс--ньютоновской компоненты
\begin{equation}
    \mathbf{G}(\mathbf{w})
    =
    \mathbf{J}(\mathbf{w})^\top
    \mathbf{A}(\mathbf{w})
    \mathbf{J}(\mathbf{w}),
\end{equation}
поэтому
\begin{equation}
    \|\mathbf{G}(\mathbf{w})\|_2
    \le
    \|\mathbf{A}(\mathbf{w})\|_2
    \|\mathbf{J}(\mathbf{w})\|_2^2
    \le
    M_{\mathbf{A}}\|\mathbf{J}(\mathbf{w})\|_2^2.
\end{equation}
Так как
\begin{equation}
    \mathbf{J}(\mathbf{w})\mathbf{J}(\mathbf{w})^\top
    =
    \sum_{i\in\{Q,K,V,1,2\}}
    \mathbf{J}^{(i)}\mathbf{J}^{(i)\top},
\end{equation}
то
\begin{equation}
    \|\mathbf{J}(\mathbf{w})\|_2^2
    =
    \|\mathbf{J}(\mathbf{w})\mathbf{J}(\mathbf{w})^\top\|_2
    \le
    \sum_{i\in\{Q,K,V,1,2\}}
    \|\mathbf{J}^{(i)}\|_2^2.
\end{equation}
Получаем \eqref{eq:app-tr-GN-general-bound}.

Оценим теперь отдельные блоки.

\textbf{1. Оценка $\mathbf{J}^{(1)}$.}
По теореме \ref{thm:app-tr-block-jacobians}
\begin{equation}
    \mathbf{J}^{(1)}
    =
    \mathbf{J}_Z
    \left(
        \mathbf{I}_T \otimes \mathbf{W}_2^\top
    \right)
    \mathbf{D}_{\sigma}(\mathbf{Y}\mathbf{W}_1)
    \left(
        \mathbf{Y}\otimes \mathbf{I}_{d_{\mathrm{ff}}}
    \right).
\end{equation}
Следовательно,
\begin{align}
    \|\mathbf{J}^{(1)}\|_2
    &\le
    \|\mathbf{J}_Z\|_2
    \|\mathbf{I}_T \otimes \mathbf{W}_2^\top\|_2
    \|\mathbf{D}_{\sigma}(\mathbf{Y}\mathbf{W}_1)\|_2
    \|\mathbf{Y}\otimes \mathbf{I}_{d_{\mathrm{ff}}}\|_2 \\
    &=
    \|\mathbf{J}_Z\|_2
    \|\mathbf{W}_2\|_2
    \cdot 1 \cdot
    \|\mathbf{Y}\|_2.
\end{align}




\textbf{2. Оценка $\mathbf{J}^{(2)}$.}
Аналогично
\begin{equation}
    \mathbf{J}^{(2)}
    =
    \mathbf{J}_Z
    \left(
        \sigma(\mathbf{Y}\mathbf{W}_1)\otimes \mathbf{I}_d
    \right),
\end{equation}
поэтому
\begin{equation}
    \|\mathbf{J}^{(2)}\|_2
    \le
    \|\mathbf{J}_Z\|_2
    \|\sigma(\mathbf{Y}\mathbf{W}_1)\|_2.
\end{equation}
Так как ReLU является $1$-липшицевой и $\sigma(\mathbf{0})=\mathbf{0}$,
\begin{equation}
    \|\sigma(\mathbf{Y}\mathbf{W}_1)\|_2
    \le
    \|\mathbf{Y}\mathbf{W}_1\|_2
    \le
    \|\mathbf{Y}\|_2\|\mathbf{W}_1\|_2.
\end{equation}




\textbf{3. Оценка $\mathbf{J}^{(k)}$, $k\in\{Q,K,V\}$.}
Из теоремы \ref{thm:app-tr-block-jacobians}
\begin{equation}
    \mathbf{J}^{(k)}
    =
    \mathbf{J}_Z \mathbf{J}_{SY}\mathbf{J}_Y \mathbf{G}_k.
\end{equation}
Следовательно,
\begin{equation}
    \|\mathbf{J}^{(k)}\|_2
    \le
    \|\mathbf{J}_Z\|_2
    \|\mathbf{J}_{SY}\|_2
    \|\mathbf{J}_Y\|_2
    \|\mathbf{G}_k\|_2.
\end{equation}
Оценим $\|\mathbf{J}_{SY}\|_2$:
\begin{align}
    \|\mathbf{J}_{SY}\|_2
    &=
    \left\|
        \mathbf{I}_{Td}
        +
        \left(
            \mathbf{I}_T \otimes \mathbf{W}_2^\top
        \right)
        \mathbf{D}_{\sigma}(\mathbf{Y}\mathbf{W}_1)
        \left(
            \mathbf{I}_T \otimes \mathbf{W}_1^\top
        \right)
    \right\|_2 \\
    &\le
    1
    +
    \|\mathbf{W}_2\|_2
    \|\mathbf{W}_1\|_2.
\end{align}
Подстановка завершает доказательство.
\end{proof}

\begin{theorem}
    \label{thm:app-tr-GN-explicit}
    Пусть выполнено
    \begin{equation}
        \|\mathbf{A}(\mathbf{w})\|_2 \le M_{\mathbf{A}}.
    \end{equation}
    Тогда
    \begin{multline}
        \label{eq:app-tr-GN-explicit}
        \|\mathbf{G}(\mathbf{w})\|_2
        \le
        M_{\mathbf{A}}
        \|\mathbf{J}_Z\|_2^2
        \Bigg[
            \|\mathbf{W}_2\|_2^2\|\mathbf{Y}\|_2^2
            +
            \|\mathbf{W}_1\|_2^2\|\mathbf{Y}\|_2^2
            +
            \\
            +
            \left(
                \|\mathbf{W}_2\|_2\|\mathbf{W}_1\|_2 + 1
            \right)^2
            \|\mathbf{J}_Y\|_2^2
            \left(
                T^2\|\mathbf{X}\|_2^2
                +
                \frac{C_{\mathrm{sm}}^2}{d_k}
                \|\mathbf{W}_V\|_2^2\|\mathbf{W}_K\|_2^2\|\mathbf{X}\|_2^6
                +
                \frac{C_{\mathrm{sm}}^2}{d_k}
                \|\mathbf{W}_V\|_2^2\|\mathbf{W}_Q\|_2^2\|\mathbf{X}\|_2^6
            \right)
        \Bigg].
    \end{multline}
\end{theorem}

\begin{proof}
По теореме \ref{thm:app-tr-GN-general-bound}
\begin{equation}
    \|\mathbf{G}(\mathbf{w})\|_2
    \le
    M_{\mathbf{A}}
    \left(
        \|\mathbf{J}^{(Q)}\|_2^2
        +
        \|\mathbf{J}^{(K)}\|_2^2
        +
        \|\mathbf{J}^{(V)}\|_2^2
        +
        \|\mathbf{J}^{(1)}\|_2^2
        +
        \|\mathbf{J}^{(2)}\|_2^2
    \right).
\end{equation}
Для FFN-параметров имеем
\begin{equation}
    \|\mathbf{J}^{(1)}\|_2^2
    \le
    \|\mathbf{J}_Z\|_2^2
    \|\mathbf{W}_2\|_2^2
    \|\mathbf{Y}\|_2^2,
\end{equation}
\begin{equation}
    \|\mathbf{J}^{(2)}\|_2^2
    \le
    \|\mathbf{J}_Z\|_2^2
    \|\mathbf{W}_1\|_2^2
    \|\mathbf{Y}\|_2^2.
\end{equation}
Для attention-параметров, используя лемму \ref{lemma:app-tr-SA-jacobian-bounds}, получаем
\begin{equation}
    \|\mathbf{J}^{(V)}\|_2^2
    \le
    \|\mathbf{J}_Z\|_2^2
    \left(
        \|\mathbf{W}_2\|_2\|\mathbf{W}_1\|_2 + 1
    \right)^2
    \|\mathbf{J}_Y\|_2^2
    T^2\|\mathbf{X}\|_2^2,
\end{equation}
\begin{equation}
    \|\mathbf{J}^{(Q)}\|_2^2
    \le
    \|\mathbf{J}_Z\|_2^2
    \left(
        \|\mathbf{W}_2\|_2\|\mathbf{W}_1\|_2 + 1
    \right)^2
    \|\mathbf{J}_Y\|_2^2
    \frac{C_{\mathrm{sm}}^2}{d_k}
    \|\mathbf{W}_V\|_2^2\|\mathbf{W}_K\|_2^2\|\mathbf{X}\|_2^6,
\end{equation}
\begin{equation}
    \|\mathbf{J}^{(K)}\|_2^2
    \le
    \|\mathbf{J}_Z\|_2^2
    \left(
        \|\mathbf{W}_2\|_2\|\mathbf{W}_1\|_2 + 1
    \right)^2
    \|\mathbf{J}_Y\|_2^2
    \frac{C_{\mathrm{sm}}^2}{d_k}
    \|\mathbf{W}_V\|_2^2\|\mathbf{W}_Q\|_2^2\|\mathbf{X}\|_2^6.
\end{equation}
Суммируя все пять слагаемых, приходим к \eqref{eq:app-tr-GN-explicit}.
\end{proof}

\begin{corollary}
    \label{cor:app-tr-GN-uniform}
    Пусть
    \begin{equation}
        \|\mathbf{A}(\mathbf{w})\|_2 \le M_{\mathbf{A}},
        \qquad
        \|\mathbf{X}\|_2 \le M_{\mathbf{X}},
        \qquad
        \|\mathbf{Y}\|_2 \le M_{\mathbf{Y}},
    \end{equation}
    \begin{equation}
        \|\mathbf{W}_Q\|_2,\|\mathbf{W}_K\|_2,\|\mathbf{W}_V\|_2,\|\mathbf{W}_1\|_2,\|\mathbf{W}_2\|_2
        \le
        M_{\mathbf{W}},
    \end{equation}
    \begin{equation}
        \sigma_{\min}(\mathbf{U}) \ge \sigma_0,
        \qquad
        \sigma_{\min}(\mathbf{S}) \ge \sigma_0 > 0.
    \end{equation}
    Обозначим
    \begin{equation}
        C_{\mathrm{LN}}(M)
        :=
        \frac{1}{\sigma_0}
        +
        \frac{M^2}{\sqrt d\,\sigma_0^3}.
    \end{equation}
    Тогда
    \begin{multline}
        \label{eq:app-tr-GN-uniform}
        \|\mathbf{G}(\mathbf{w})\|_2
        \le
        M_{\mathbf{A}}
        C_{\mathrm{LN}}(\|\mathbf{S}\|_2)^2
        \Bigg[
            2M_{\mathbf{Y}}^2M_{\mathbf{W}}^2
            +
            \\
            +
            \left(
                M_{\mathbf{W}}^2 + 1
            \right)^2
            C_{\mathrm{LN}}(\|\mathbf{U}\|_2)^2
            \left(
                T^2 M_{\mathbf{X}}^2
                +
                \frac{2C_{\mathrm{sm}}^2}{d_k}
                M_{\mathbf{W}}^4 M_{\mathbf{X}}^6
            \right)
        \Bigg].
    \end{multline}
\end{corollary}

\begin{proof}
По лемме \ref{lemma:app-tr-LN-jacobian-bound}
\begin{equation}
    \|\mathbf{J}_Y\|_2
    \le
    C_{\mathrm{LN}}(\|\mathbf{U}\|_2),
    \qquad
    \|\mathbf{J}_Z\|_2
    \le
    C_{\mathrm{LN}}(\|\mathbf{S}\|_2).
\end{equation}
Далее в теореме \ref{thm:app-tr-GN-explicit} подставляем равномерные оценки
\begin{equation}
    \|\mathbf{W}_Q\|_2,\|\mathbf{W}_K\|_2,\|\mathbf{W}_V\|_2,\|\mathbf{W}_1\|_2,\|\mathbf{W}_2\|_2
    \le M_{\mathbf{W}},
\end{equation}
\begin{equation}
    \|\mathbf{X}\|_2 \le M_{\mathbf{X}},
    \qquad
    \|\mathbf{Y}\|_2 \le M_{\mathbf{Y}}.
\end{equation}
Первые два члена дают
\begin{equation}
    \|\mathbf{W}_2\|_2^2\|\mathbf{Y}\|_2^2
    +
    \|\mathbf{W}_1\|_2^2\|\mathbf{Y}\|_2^2
    \le
    2M_{\mathbf{Y}}^2M_{\mathbf{W}}^2.
\end{equation}
Для attention-части получаем
\begin{equation}
    \|\mathbf{W}_2\|_2\|\mathbf{W}_1\|_2 + 1
    \le
    M_{\mathbf{W}}^2 + 1,
\end{equation}
а сумма двух старших членов порядка $\|\mathbf{X}\|_2^6$ оценивается как
\begin{equation}
    \frac{C_{\mathrm{sm}}^2}{d_k}
    \|\mathbf{W}_V\|_2^2\|\mathbf{W}_K\|_2^2\|\mathbf{X}\|_2^6
    +
    \frac{C_{\mathrm{sm}}^2}{d_k}
    \|\mathbf{W}_V\|_2^2\|\mathbf{W}_Q\|_2^2\|\mathbf{X}\|_2^6
    \le
    \frac{2C_{\mathrm{sm}}^2}{d_k}M_{\mathbf{W}}^4M_{\mathbf{X}}^6.
\end{equation}
Подстановка этих оценок в \eqref{eq:app-tr-GN-explicit} завершает доказательство.
\end{proof}